\section{Integral comparison with constant periods along an inert arc}

% BEGIN J82F1
\paragraph{J82F1. Two integral matrices at the full backward modulus.}
Fix either tested crystalline lattice on an arc of I82B3 ending
at $c$, of length $\ell\ge j$. Write $i=c-\ell$, $A=O[[v]]$,
$L=O[1/p]$, and let $e$ be a row of basis vectors of the tame
invariant Wach lattice $N_i$. Put $q=p-1$ and use the parameters
$M,g,\sigma$, factors $P_s$, and $f_j$ of I82B4. The
Frobenius span at $c$ is free on the row $E=\varphi^\ell(e)$:
its columns generate by definition and are independent in the
\'etale module by invertibility of linearized Frobenius there.
The relevant modulus at $c$ is
\[
  f'=\varphi^{\ell-j}(f_j)
  \ \sim\ v^M\!\prod_{0\le s<\ell-j}P_s^{qM}
                 \prod_{\ell-j\le s<\ell}P_s^{3g},
\]
where $\sim$ denotes association in $A$. This follows by
iteration of $\varphi(v)\sim vP_0^q$. Put $d_s=qM$ in the
first range and $d_s=3g$ in the second. Thus $f'$ is a unit
times a distinguished monic polynomial and $A/(f')$ is
$p$-torsion-free.

Write $\gamma(e)=eG(v)$ and solve $G\gamma(U)=U$ in
$\mathrm{GL}_r(L[[v]])$ with $U(0)=1$. The equation in
positive degree $k$ divides an integral combination of previous
coefficients by $1-(1+p)^{qk}$. Its valuation is $1+v_p(k)$
and $G(0)=1$. Induction gives the denominator bound
$p^{w(k)}$, $w(k)=k+v_p(k!)$, in degree $k$.
The inverse satisfies $\gamma(U^{-1})=U^{-1}G$, giving the
same bound by the same recursion. This is the normalized
invariant frame of L78W9 in tame coordinates. In particular
the polynomial truncations of $U,U^{-1}$ below degree $M$
are inverse modulo $v^M$ over $L$.

For each of these polynomials I82B4 gives an integral
representative after applying $p^\sigma\varphi^j$ modulo $f_j$,
by its estimate on every monomial of denominator $p^{w(k)}$.
Apply $\varphi^{\ell-j}$ to the integral representatives.
This gives matrices $B_+,B_-$ over $A$ representing
$p^\sigma\varphi^\ell(U\bmod v^M)$ and
$p^\sigma\varphi^\ell(U^{-1}\bmod v^M)$ respectively
over $(A/(f'))[1/p]$. Since $f_j\mid\varphi^j(v)^M$,
$f'\mid\varphi^\ell(v)^M$. Hence both products
$B_+B_-$ and $B_-B_+$ equal $p^{2\sigma}I$ modulo $f'$,
first rationally and then integrally by the torsion-freeness
above. This one pair of matrices serves at every position
of the backward arc.
% END J82F1

% BEGIN J82F2
\paragraph{J82F2. The two comparisons at each primitive expansion.}
Let $D_i$ and $D_c$ be the crystalline components of the two
letters. The constant period represented near $v=0$ by $eU$ is
the row $d=\mathrm{cmp}(e)$ in $D_i$. Indeed in the tied period
basis of L78W9--L78W11 the columns of $eU$ are invariant, and a
generator-invariant scalar series there is constant, by the
first nonzero coefficient test; its value at zero is determined
by $U(0)=1$. The comparison at zero and the invertible Frobenius
of W79H3 now identify the native constant basis at $c$ with
$\lambda=\varphi_{\rm cris}^\ell d$. For each $0\le s<\ell$,
take every primitive root $\zeta$ of order $p^{s+1}$ in a common
splitting coefficient field and expand at $1+X=\zeta\exp(Y)$.
After $\ell$ Frobenius steps this is the expansion at
$1+X=\exp(p^\ell Y)$ in letter $i$, with crystalline
Frobenius applied to its vector coefficients. In particular
the Taylor matrix of $E$ is regular with constant row
$\lambda$, a basis. In these expansions
$E\varphi^\ell(U)=\lambda$ as a formal identity: the argument
of $U$ is now $\varphi^\ell(v)$, which vanishes to order $q$
because $v$ has a zero of order $q$ at $X=0$.
Thus the finite truncations in J82F1 suffice to compare the
two bases through all degrees below $d_s\le qM$.

Set $\Lambda=\prod_{s<\ell}P_s^{H_s}$ with the original
arc exponents. The physical numerator span is $EA^r$, and the
physical module $W$ in the partner condition is $\Lambda^{-1}EA^r$.
Identify constant numerators with columns $a\in A^r$, whose
projection at a root is $\lambda a/\Lambda$. All the $P_t$
for $t\ne s$ and $v$ are units at this root. The factor $P_s$
has a simple zero, since it is associated in the original
variable to the primitive cyclotomic factor, which is separable.
Consequently the error space $\mathrm{Fl}_s$ consists of vectors
of $Y$-order at least $d_s-H_s$: its defining test on the
numerator is divisibility by the expansion of $f'$.
For $a\in A^r$ take the physical vector $EB_+a/\Lambda$.
Its Taylor series agrees with $p^\sigma\lambda a/\Lambda$
modulo $\mathrm{Fl}_s$, by J82F1 and regularity of $E$.
Conversely, for $Eb/\Lambda$ with $b\in A^r$, the integral
constant numerator $B_-b$ agrees in the same sense with
$p^\sigma Eb/\Lambda$, by the inverse frame identity.
These are both clauses of $\mathrm{Partner}(\sigma)$, with
one representative for all primitive roots and all levels.
% END J82F2

% BEGIN J82F3
\paragraph{J82F3. Localization of almost invariant vectors and functoriality.}
Let $x\in N_i$ satisfy $(\gamma-1)x\in v^MN_i$ and
$x\equiv ea\bmod vN_i$ for $a\in O^r$. Write $x=eU b(v)$
formally at zero. Then $b(0)=a$ and
$\gamma b-b\in v^ML[[v]]^r$ since $eU$ is invariant.
If $b-a$ had first nonzero degree $0<k<M$, the same coefficient
of $\gamma b-b$ would be multiplied by $(1+p)^{qk}-1\ne0$.
Thus $b\equiv a\bmod v^M$. For the physical vector
$w=\varphi^\ell(x)/\Lambda$, J82F2 now gives
\[
    \Lambda\,\mathrm{loc}_s(w)\equiv\lambda a\pmod{Y^{qM}}
             \qquad(0\le s<\ell)
\]
in the projected Taylor coordinates at every root.
Here $qM>h+2g\ge3g$ by I82B4, so $qM\ge d_s$ throughout
the arc. After division by $\Lambda$ this is precisely
$\mathrm{loc}_s(w)-\mathrm{locL}_s(a)\in\mathrm{Fl}_s$,
the native $\mathrm{LocalizesConst}(M,e)$.

If $f(e_1)=e_2 F(v)$ is the tested map at $i$, the coordinates
$U_2^{-1}F(v)U_1$ are generator-invariant and therefore constant,
with value $F(0)$. Applying $\varphi^\ell$ shows in the
crystalline bases that the map on constant numerators is
coefficientwise $F(0)$. Truncating the same identity as in
J82F1 also gives compatibility of each scaled comparison modulo
$f'$. These statements supply the physical frame, its error
order, and its naturality for both tested lattices. Lifting
a vector with its prescribed filtered residue under this map
is a further condition in H82A3.
% END J82F3
